\documentclass[11pt,twocolumn]{article}

\usepackage[margin=0.75in,top=0.85in,bottom=0.85in]{geometry}
\usepackage{times}
\usepackage{microtype}
\usepackage{booktabs}
\usepackage{multirow}
\usepackage{graphicx}
\usepackage{amsmath}
\usepackage{hyperref}
\usepackage{xcolor}
\usepackage{enumitem}
\usepackage{caption}
\usepackage[numbers,compress]{natbib}

\hypersetup{colorlinks=true,linkcolor=black,citecolor=black,urlcolor=blue}
\setlength{\columnsep}{0.25in}

% ─── compact section spacing ──────────────────────────────────────────────────
\usepackage{titlesec}
\titlespacing*{\section}{0pt}{6pt}{3pt}
\titlespacing*{\subsection}{0pt}{4pt}{2pt}
\setlength{\parskip}{2pt}

\title{\textbf{Conversational Hallucination Drift: An Episodic Retrieval\\
Framework and Error Taxonomy for Long-Term Memory Evaluation}}

\author{
  Kamesh R. \\
  \small{\texttt{kamesh@episodic-log}}
}

\date{}

\begin{document}

\maketitle
\thispagestyle{empty}

% ─── ABSTRACT ────────────────────────────────────────────────────────────────
\begin{abstract}
Long-term conversational memory remains a persistent weakness of large
language models. We study this failure mode on LongMemEval~\cite{wu2025longmemeval},
a 500-question benchmark embedded in realistic user-assistant histories.
We introduce the \textbf{Conversational Hallucination Drift (CHD)} taxonomy, which
classifies memory errors into four mechanistically distinct types:
\textit{commission}, \textit{omission}, \textit{distortion}, and
\textit{confabulation}. We implement an open episodic retrieval
framework pairing three summarization strategies (lexical, scout,
echo) with two retrieval conditions (recall, grep-recall) and evaluate
Qwen3-32B as the reasoning model. Our best configuration---recall
with scout summaries---achieves \textbf{24.9\% accuracy}, closely matching
the 24.7\% reported for Coze with GPT-3.5-turbo~\cite{wu2025longmemeval},
using a fully open-weight stack. CHD analysis reveals that \textbf{omission
accounts for 68\% of remaining errors}, identifying generation
faithfulness---not retrieval---as the primary bottleneck. All code
and evaluation scripts are released.
\end{abstract}

% ─── 1. INTRODUCTION ─────────────────────────────────────────────────────────
\section{Introduction}

Conversational agents are increasingly expected to remember what users
told them days or weeks ago. Yet evaluations consistently show that
LLMs degrade sharply when answer evidence is buried in long interaction
histories. LongMemEval~\cite{wu2025longmemeval} operationalises this
challenge: 500 questions require reasoning over sessions that can
exceed 115k tokens, and the best open-model long-context baseline
(GPT-4o, full context) achieves only 60.6\%, while smaller open
models fall to 12--33\%.

Memory-augmented systems address this with retrieval scaffolds, but
most published work uses proprietary APIs, leaving the open-weight
community without a reproducible reference point. Separately, existing
hallucination taxonomies~\cite{li2023halueval,maynez2020faithfulness}
are designed for single-turn factuality and do not capture the
\emph{self-referential} errors that arise when a model misremembers
its own past behaviour.

We make three contributions:
\begin{enumerate}[leftmargin=*,itemsep=0pt,topsep=2pt]
  \item A \textbf{CHD taxonomy} (§\ref{sec:taxonomy}) grounded in
        cognitive-science confabulation research~\cite{schnider2008confabulating}
        that distinguishes four mechanistically distinct error types.
  \item An \textbf{open episodic retrieval framework} (§\ref{sec:framework})
        that ingests LongMemEval sessions and evaluates six condition
        combinations in a fully reproducible pipeline.
  \item A \textbf{controlled ablation} (§\ref{sec:experiments}) showing
        that recall with scout summarisation reaches 24.9\% with
        Qwen3-32B, closely matching the commercial Coze/GPT-3.5 baseline,
        and that omission---not retrieval failure---accounts for
        the majority of remaining errors.
\end{enumerate}

% ─── 2. CHD TAXONOMY ─────────────────────────────────────────────────────────
\section{CHD Taxonomy}
\label{sec:taxonomy}

Existing hallucination taxonomies distinguish \emph{intrinsic} errors
(contradicting source material) from \emph{extrinsic} ones (inventing
unsupported content)~\cite{maynez2020faithfulness}. This two-way split
is too coarse for memory agents, where the ``source'' is the model's own
interaction history. We propose four types:

\begin{itemize}[leftmargin=*,itemsep=0pt,topsep=2pt]
  \item \textbf{Commission.} The model asserts it performed an action it
        never did (\emph{false memory insertion}).
  \item \textbf{Omission.} The model fails to recall a clearly stated past
        action or fact (\emph{memory dropout}).
  \item \textbf{Distortion.} The model recalls the correct event but with
        wrong key details---wrong date, wrong value, wrong name.
  \item \textbf{Confabulation.} The model invents plausible-sounding history
        with no grounding in the log.
\end{itemize}

A judge model (Qwen2.5-14B-Instruct~\cite{qwen25}) classifies each
predicted answer into one of these four categories or \textit{correct},
given the question, ground truth, and predicted answer. We note that
our judge differs from the GPT-4o judge used in the original
LongMemEval paper; direct accuracy comparisons should be interpreted
with this caveat in mind. Inter-annotator agreement with human labels
is deferred to future work.

% ─── 3. EPISODIC RETRIEVAL FRAMEWORK ─────────────────────────────────────────
\section{Episodic Retrieval Framework}
\label{sec:framework}

\paragraph{Ingestion.} Each LongMemEval session is converted into an
append-only \texttt{log.jsonl} file (one \texttt{TurnEvent} per line),
preserving role, content, and turn ID.

\paragraph{Summarisation.} Three strategies build a per-session summary
index:
\begin{itemize}[leftmargin=*,itemsep=0pt,topsep=2pt]
  \item \textbf{Lexical.} Rule-based keyword extraction; no model; runs on
        CPU in seconds. Produces sparse, reproducible summaries.
  \item \textbf{Scout.} Qwen2.5-7B-Instruct~\cite{qwen25} writes a
        retrieval-optimised 1--2 sentence summary per turn, emphasising
        specific nouns, dates, and action verbs.
  \item \textbf{Echo.} The same 7B model writes a detail-preserving
        summary, retaining numbers, names, and quoted values verbatim.
\end{itemize}
Each summary is capped at 120 characters to keep the full 500-turn
index within 25k tokens.

\paragraph{Retrieval conditions.} We implement two agent-loop strategies:
\begin{itemize}[leftmargin=*,itemsep=0pt,topsep=2pt]
  \item \textbf{Recall.} The full summary index is injected into the
        first user message. The model scans it, identifies relevant
        turn IDs, and calls \texttt{load\_turn} to read verbatim content
        before answering. Up to 8 \texttt{load\_turn} calls permitted.
  \item \textbf{Grep-recall.} No index is injected. The model formulates
        keywords and calls \texttt{grep\_memory} (substring match over
        summaries); the top-3 matches are returned with full verbatim
        content. Up to 15 combined tool calls permitted.
\end{itemize}
Both strategies use Qwen3-32B~\cite{qwen3} via vLLM~\cite{kwon2023vllm}
with tensor-parallel inference across 8 A100 40\,GB GPUs.

% ─── 4. EXPERIMENTS ──────────────────────────────────────────────────────────
\section{Experiments}
\label{sec:experiments}

\subsection{Setup}

We evaluate all 500 questions from LongMemEval~\cite{wu2025longmemeval}.
The reasoning model is Qwen3-32B (BF16, tp=8). The summariser for
scout and echo conditions is Qwen2.5-7B-Instruct (BF16, tp=4). Judging
uses Qwen2.5-14B-Instruct (BF16, tp=8). All inference runs on 8$\times$
A100 40\,GB via vLLM. The \textit{amnesiac} baseline answers purely from
parametric knowledge with no memory access.

\subsection{Main Results}

Table~\ref{tab:main} reports accuracy and CHD breakdown for all
conditions. Published baselines from~\citet{wu2025longmemeval} are
included for reference (marked $\dagger$; GPT-4o judge).

\begin{table}[h]
\centering
\caption{Accuracy and CHD error breakdown on LongMemEval (500 Qs).
$\dagger$ = published result with GPT-4o judge; all others use
Qwen2.5-14B judge. Com.~= commission, Omi.~= omission,
Dis.~= distortion, Con.~= confabulation.}
\label{tab:main}
\resizebox{\columnwidth}{!}{%
\begin{tabular}{llccccc}
\toprule
\textbf{Condition} & \textbf{Method} & \textbf{Acc.} & \textbf{Com.} & \textbf{Omi.} & \textbf{Dis.} & \textbf{Con.} \\
\midrule
\multicolumn{2}{l}{\textit{Published baselines}$^\dagger$} & & & & & \\
GPT-4o full ctx & --- & 60.6\% & --- & --- & --- & --- \\
ChatGPT memory & GPT-4o & 57.7\% & --- & --- & --- & --- \\
Coze memory & GPT-3.5 & 24.7\% & --- & --- & --- & --- \\
\midrule
\multicolumn{2}{l}{\textit{Our conditions (Qwen3-32B)}} & & & & & \\
recall      & scout   & \textbf{24.9\%} & 2.8\% & 53.4\% & 16.4\% & 4.4\% \\
grep-recall & scout   & 22.2\% & 2.4\% & 49.4\% & 19.2\% & 6.8\% \\
grep-recall & lexical & 20.4\% & 2.2\% & 59.8\% & 14.2\% & 3.4\% \\
recall      & lexical & 19.8\% & 2.2\% & 59.2\% & 13.8\% & 5.0\% \\
recall      & echo    & 13.2\% & 3.4\% & 62.0\% & 14.0\% & 7.4\% \\
grep-recall & echo    & 12.6\% & 0.8\% & 68.6\% & 13.2\% & 4.8\% \\
amnesiac    & ---     &  5.4\% & 0.6\% & 80.4\% &  7.4\% & 6.2\% \\
\bottomrule
\end{tabular}
}
\end{table}

\subsection{Analysis}

\paragraph{Retrieval lifts performance substantially.}
The amnesiac baseline reaches only 5.4\%, relying entirely on
parametric knowledge. Our best retrieval condition (recall/scout)
reaches 24.9\%---a 19.5 percentage-point improvement, confirming
that memory scaffolding provides a substantial and measurable signal.

\paragraph{Scout summarisation consistently outperforms lexical and echo.}
Recall/scout (24.9\%) exceeds recall/lexical (19.8\%) by 5.1\,pp and
recall/echo (13.2\%) by 11.7\,pp. Scout's retrieval-optimised phrasing
allows the model to quickly identify relevant turns from the summary
index. Echo summaries, despite being more information-dense, appear to
introduce noise: the model attends to verbose detail in the index
rather than calling \texttt{load\_turn} to read authoritative verbatim
content.

\paragraph{Grep-recall nearly matches recall at equal tool budgets.}
With 15 tool calls, grep-recall/scout reaches 22.2\% versus
recall/scout at 24.9\%---a gap of only 2.7\,pp. Earlier runs with
grep-recall capped at 5 tool calls (versus recall's 8) showed a
larger gap, confirming that the budget asymmetry, not the retrieval
strategy itself, drove the apparent advantage of recall. Grep-recall
is preferable in settings where injecting the full summary index
is impractical due to context constraints.

\paragraph{Omission dominates all conditions.}
Across all memory conditions, omission is the leading error type,
accounting for 49--68\% of failures (Table~\ref{tab:main}). Commission
and confabulation are consistently below 8\%, indicating the model
rarely invents memories; instead, it retrieves correctly but produces
answers that drop precise values---a date becomes ``Tuesday'' without
the month, an amount loses its exact figure. This diagnosis
distinguishes two possible improvement axes: retrieval quality
(already substantially addressed) and \emph{generation faithfulness}
(the remaining bottleneck).

\paragraph{Comparison to commercial baselines.}
Recall/scout achieves 24.9\%, marginally exceeding the 24.7\% reported
for Coze with GPT-3.5-turbo~\cite{wu2025longmemeval}. We emphasise that
this comparison is approximate: the original paper uses GPT-4o as
judge, while we use Qwen2.5-14B. Nevertheless, the results suggest
that a lightweight open-weight pipeline can achieve competitive
accuracy relative to commercial memory systems at this performance
tier, while providing full interpretability via CHD error breakdown.
The gap to frontier systems (ChatGPT/GPT-4o at 57.7\%) remains large
and is not addressed by our current framework.

% ─── 5. RELATED WORK ─────────────────────────────────────────────────────────
\section{Related Work}

\paragraph{Long-term memory benchmarks.}
LongMemEval~\cite{wu2025longmemeval} is the primary benchmark we use.
LOCOMO~\cite{maharana2024locomo} is a complementary benchmark used
by most recent memory-system papers; we leave cross-benchmark
evaluation to future work.

\paragraph{Memory-augmented agents.}
MemGPT~\cite{packer2023memgpt} introduces a paged virtual context for
LLM agents, analogous to OS memory management. Recent systems---A-MEM~\cite{xu2025amem}, Mem0~\cite{chhikara2025mem0}, and
MemoryOS~\cite{kang2025memoryos}---add dynamic memory update and
hierarchical storage, evaluated primarily on LOCOMO rather than
LongMemEval. ReadAgent~\cite{lee2024readagent} compresses documents
into gist memories and retrieves originals on demand, closely related
to our recall condition with scout summaries.

\paragraph{Hallucination taxonomies.}
Intrinsic/extrinsic faithfulness~\cite{maynez2020faithfulness} and
HaluEval~\cite{li2023halueval} characterise factuality errors in
summarisation and QA. Our CHD taxonomy focuses on the self-referential
domain: errors arising when a model misremembers its own past
interactions, motivated by cognitive-science accounts of
confabulation~\cite{schnider2008confabulating}.

% ─── 6. CONCLUSION ───────────────────────────────────────────────────────────
\section{Conclusion}

We introduced the CHD error taxonomy and an open episodic retrieval
framework, showing that recall with scout summarisation and
Qwen3-32B achieves 24.9\% on LongMemEval---competitive with a
commercial GPT-3.5 memory system, with the caveat that judge models
differ. CHD analysis reveals that omission accounts for the majority
of remaining failures, pointing to generation faithfulness as the
primary target for future improvement. The gap to frontier systems
(GPT-4o) remains large and underscores that the problem is far from
solved. Our code, evaluation pipeline, and prompts are publicly
released to facilitate reproducible follow-up.

\textbf{Limitations.} We evaluate a single model size (32B) on a single
benchmark. Judge model mismatch limits strict numerical comparability
with published baselines. CHD inter-annotator agreement with human
labels has not been measured. We intend to address these in follow-up
work.

% ─── REFERENCES ──────────────────────────────────────────────────────────────
\bibliographystyle{abbrvnat}
\bibliography{references}

\end{document}
